\sectionkicker{Source-backed technical notes}
\subsection*{Multimodal \& Generation — 独立して厚くなる生成・理解・grounding: Technical Notes}
本欄は記事本文で圧縮した一次資料上の情報を、比較・再検証しやすい形で整理したものである。時系列、確認済みの主張、留意点、一次資料URLを掲載する。
\medskip
\subsection*{Theme at a glance}
\addcontentsline{toc}{subsection}{Theme at a glance}
\begin{center}
\footnotesize
\begin{tabularx}{\linewidth}{@{}p{0.20\linewidth}p{0.10\linewidth}p{0.14\linewidth}X@{}}
\toprule
資料 & 位置づけ & 種別 & 時系列 \\
\midrule
LongCat-Next: Lexicalizing Modalities as Discrete Tokens & 主要資料 & 論文 & 2026-03-29 (論文公開) \\
Unify-Agent: A Unified Multimodal Agent for World-Grounded Image Synthesis & 主要資料 & 論文 & 2026-03-31 (論文公開) \\
DreamLite: A Lightweight On-Device Unified Model for Image Generation and Editing & 補足資料 & 論文 & 2026-03-30 (論文公開) \\
Gen-Searcher: Reinforcing Agentic Search for Image Generation & 補足資料 & 論文 & 2026-03-30 (論文公開) \\
\bottomrule
\end{tabularx}
\end{center}
\normalsize

\begin{technicalnote}{LongCat-Next: Lexicalizing Modalities as Discrete Tokens}{主要資料}
\begingroup
% reader-facing Technical Notes break policy
\widowpenalty=10000
\clubpenalty=10000
\displaywidowpenalty=10000
\begin{tabularx}{\linewidth}{@{}>{\bfseries}p{0.22\linewidth}X@{}}
組織 & - \\
種別 & 論文 \\
時系列 & 2026-03-29 (論文公開) \\
\end{tabularx}
\smallskip
{\bfseries 一次資料から整理したtechnical points}
\begin{itemize}[leftmargin=1.5em,itemsep=0.35em]
\item \textbf{Author claim}: 著者らは、multimodal情報を共通の離散空間で表現し、modalitiesを横断して一貫したautoregressive modelingを行うDiNA（Discrete Native Autoregressive）を提案している。
\end{itemize}
\begin{minipage}{\linewidth}
% reader-facing Technical Notes generic boundary/source tail group
{\bfseries 読む際の境界}
\begin{itemize}[leftmargin=1.5em,itemsep=0.35em]
\item \textbf{分析上の留意点}: 論文中の結果・ベンチマーク値は著者報告であり、本サーベイでは独立再現していない。
\end{itemize}
\begin{samepage}
{\bfseries 一次資料}
\begingroup\sloppy
\begin{itemize}[leftmargin=1.5em,itemsep=0.25em]
\item {\scriptsize\url{http://arxiv.org/abs/2603.27538v1}}
\end{itemize}
\endgroup
\end{samepage}
\end{minipage}
\endgroup
\end{technicalnote}
\medskip

\begin{technicalnote}{Unify-Agent: A Unified Multimodal Agent for World-Grounded Image Synthesis}{主要資料}
\begingroup
% reader-facing Technical Notes break policy
\widowpenalty=10000
\clubpenalty=10000
\displaywidowpenalty=10000
\begin{tabularx}{\linewidth}{@{}>{\bfseries}p{0.22\linewidth}X@{}}
組織 & - \\
種別 & 論文 \\
時系列 & 2026-03-31 (論文公開) \\
\end{tabularx}
\smallskip
{\bfseries 一次資料から整理したtechnical points}
\begin{itemize}[leftmargin=1.5em,itemsep=0.35em]
\item \textbf{Author claim}: 著者らは、image generationをprompt understanding、multimodal evidence search、grounded recaptioning、final synthesisからなるagentic pipelineとして扱うUnify-Agentを提案している。
\end{itemize}
\begin{minipage}{\linewidth}
% reader-facing Technical Notes generic boundary/source tail group
{\bfseries 読む際の境界}
\begin{itemize}[leftmargin=1.5em,itemsep=0.35em]
\item \textbf{分析上の留意点}: 論文中の結果・ベンチマーク値は著者報告であり、本サーベイでは独立再現していない。
\item \textbf{分析上の留意点}: 参照先はarXiv v2だが、3月の時系列は初回公開日を基準とする。後日の改訂内容を3月時点へ遡及して扱わない。
\end{itemize}
\begin{samepage}
{\bfseries 一次資料}
\begingroup\sloppy
\begin{itemize}[leftmargin=1.5em,itemsep=0.25em]
\item {\scriptsize\url{http://arxiv.org/abs/2603.29620v2}}
\end{itemize}
\endgroup
\end{samepage}
\end{minipage}
\endgroup
\end{technicalnote}
\medskip

\begin{technicalnote}{DreamLite: A Lightweight On-Device Unified Model for Image Generation and Editing}{補足資料}
\begingroup
% reader-facing Technical Notes break policy
\widowpenalty=10000
\clubpenalty=10000
\displaywidowpenalty=10000
\begin{tabularx}{\linewidth}{@{}>{\bfseries}p{0.22\linewidth}X@{}}
組織 & - \\
種別 & 論文 \\
時系列 & 2026-03-30 (論文公開) \\
\end{tabularx}
\smallskip
{\bfseries 一次資料から整理したtechnical points}
\begin{itemize}[leftmargin=1.5em,itemsep=0.35em]
\item \textbf{Author claim}: 著者らは、0.39B規模でT2I generationとtext-guided image editingを単一networkで扱うon-device unified diffusion model、DreamLiteを提案している。
\end{itemize}
\begin{minipage}{\linewidth}
% reader-facing Technical Notes generic boundary/source tail group
{\bfseries 読む際の境界}
\begin{itemize}[leftmargin=1.5em,itemsep=0.35em]
\item \textbf{分析上の留意点}: 論文中の結果・ベンチマーク値は著者報告であり、本サーベイでは独立再現していない。
\end{itemize}
\begin{samepage}
{\bfseries 一次資料}
\begingroup\sloppy
\begin{itemize}[leftmargin=1.5em,itemsep=0.25em]
\item {\scriptsize\url{http://arxiv.org/abs/2603.28713v1}}
\end{itemize}
\endgroup
\end{samepage}
\end{minipage}
\endgroup
\end{technicalnote}
\medskip

\begin{technicalnote}{Gen-Searcher: Reinforcing Agentic Search for Image Generation}{補足資料}
\begingroup
% reader-facing Technical Notes break policy
\widowpenalty=10000
\clubpenalty=10000
\displaywidowpenalty=10000
\begin{tabularx}{\linewidth}{@{}>{\bfseries}p{0.22\linewidth}X@{}}
組織 & - \\
種別 & 論文 \\
時系列 & 2026-03-30 (論文公開) \\
\end{tabularx}
\smallskip
{\bfseries 一次資料から整理したtechnical points}
\begin{itemize}[leftmargin=1.5em,itemsep=0.35em]
\item \textbf{Author claim}: 著者らは、multi-hop reasoningとsearchで必要なtextual knowledge・reference imageを集め、grounded generationへつなぐsearch-augmented image generation agent、Gen-Searcherを提案している。
\end{itemize}
\begin{minipage}{\linewidth}
% reader-facing Technical Notes generic boundary/source tail group
{\bfseries 読む際の境界}
\begin{itemize}[leftmargin=1.5em,itemsep=0.35em]
\item \textbf{分析上の留意点}: 論文中の結果・ベンチマーク値は著者報告であり、本サーベイでは独立再現していない。
\item \textbf{分析上の留意点}: 参照先はarXiv v3だが、3月の時系列は初回公開日を基準とする。後日の改訂内容を3月時点へ遡及して扱わない。
\end{itemize}
\begin{samepage}
{\bfseries 一次資料}
\begingroup\sloppy
\begin{itemize}[leftmargin=1.5em,itemsep=0.25em]
\item {\scriptsize\url{http://arxiv.org/abs/2603.28767v3}}
\end{itemize}
\endgroup
\end{samepage}
\end{minipage}
\endgroup
\end{technicalnote}
\medskip
